% Abstract and index terms (issue #5).

\begin{abstract}
Matrix multiplication dominates the run time of neural-network inference and many signal-processing workloads, and on general-purpose processors its cost is set by data movement rather than arithmetic. This paper presents the Phase~1 prototype of a matrix-multiplication accelerator built on an output-stationary systolic array. An $N \times N$ grid of processing elements, $4 \times 4$ by default, multiplies signed 8-bit fixed-point operands and accumulates each element of the result in place in a 32-bit register. Operands enter through two skew buffers, so each row of $A$ and each column of $B$ reaches the array one cycle later than the previous one. A five-state controller sequences the computation. A host loads $A$ and $B$ through write ports, pulses \texttt{start} and reads $C$ after \texttt{done}. The design is synthesizable, parameterized Verilog-2005 with a written timing contract: a product takes $K + 2N - 2$ compute cycles, 10 for a $4 \times 4$ array, and 12 cycles from \texttt{start} to \texttt{done}. That is 16 multiply-accumulates per cycle at peak, and 6.4 times fewer compute cycles than a single sequential multiply-accumulate unit. Every module is checked by a self-checking testbench that runs in continuous integration, and a Python golden model generates 106 directed and random test cases for the complete accelerator.
\end{abstract}

\begin{IEEEkeywords}
systolic array, matrix multiplication, hardware accelerator, processing element, output-stationary dataflow, fixed-point arithmetic, Verilog
\end{IEEEkeywords}
